%! The Inverse and Determinant of a Matrix — Unit 4, Lesson 4.11 — Group Activity (Answer Key)
\documentclass[10pt]{article}
\usepackage{precalculus-article}
\usepackage{precalculus-key}   % pulls in -boxes; do NOT also load -boxes
\newcommand{\TallMath}[1]{$\displaystyle #1\rule[-1.4em]{0pt}{3.2em}$}

\begin{document}

\pageheader{Unit 4, Lesson 4.11}{Group Activity --- Answer Key}
\namepartnerperiod

\vspace{0.12in}

\begin{tcolorbox}[colback=white, colframe=black!40,
    title=\textbf{Tier R --- Remediate},
    fonttitle=\bfseries, arc=2mm, left=3mm, right=3mm, top=2mm, bottom=2mm]
Let $A = \begin{bmatrix}3&2\\1&4\end{bmatrix}$.

\begin{enumerate}[leftmargin=1.5em, itemsep=8pt]
    \item Identify $a$, $b$, $c$, $d$ from $A$.
          \quad $a = $ \ans{3}, \; $b = $ \ans{2},
          \; $c = $ \ans{1}, \; $d = $ \ans{4}

    \item Compute $\det(A) = ad - bc$.\\[4pt]
          \ans{$\det(A) = 3(4) - 2(1) = 12 - 2 = 10$}

    \item Is $A$ invertible? Circle one and explain: \quad \textbf{Yes} \quad / \quad \textbf{No}\\[4pt]
          \ans{Yes. Since $\det(A) = 10 \ne 0$, the inverse exists.}
\end{enumerate}
\end{tcolorbox}

\vspace{0.1in}

\begin{tcolorbox}[colback=white, colframe=black!40,
    title=\textbf{Tier A --- Approaching Proficiency},
    fonttitle=\bfseries, arc=2mm, left=3mm, right=3mm, top=2mm, bottom=2mm]
Let $A = \begin{bmatrix}3&2\\1&4\end{bmatrix}$.

\begin{enumerate}[leftmargin=1.5em, itemsep=8pt]
    \item Compute $\det(A)$ and write the inverse formula for this matrix.\\[4pt]
          \ans{$\det(A) = 3(4)-2(1) = 10$. \quad
          $A^{-1} = \dfrac{1}{10}\begin{bmatrix}4&-2\\-1&3\end{bmatrix}$}

    \item Find $A^{-1}$.
          \[
              A^{-1} = \ans{$\dfrac{1}{10}$}\begin{bmatrix} \ans{4} & \ans{$-2$} \\[10pt]
              \ans{$-1$} & \ans{3} \end{bmatrix}
              = \ans{$\begin{bmatrix}0.4 & -0.2\\ -0.1 & 0.3\end{bmatrix}$}
          \]

    \item Verify your answer by computing $A \cdot A^{-1}$. Show all work.\\[4pt]
          \ans{$\begin{bmatrix}3&2\\1&4\end{bmatrix}\dfrac{1}{10}\begin{bmatrix}4&-2\\-1&3\end{bmatrix}
          = \dfrac{1}{10}\begin{bmatrix}12-2 & -6+6\\ 4-4 & -2+12\end{bmatrix}
          = \dfrac{1}{10}\begin{bmatrix}10&0\\0&10\end{bmatrix} = \begin{bmatrix}1&0\\0&1\end{bmatrix}
          = I$ \checkmark}
\end{enumerate}
\end{tcolorbox}

\vspace{0.1in}

\begin{tcolorbox}[colback=white, colframe=black!40,
    title=\textbf{Tier E --- Extension},
    fonttitle=\bfseries, arc=2mm, left=3mm, right=3mm, top=2mm, bottom=2mm]
Let $B = \begin{bmatrix}3&6\\1&2\end{bmatrix}$.

\begin{enumerate}[leftmargin=1.5em, itemsep=8pt]
    \item Compute $\det(B)$.\\[4pt]
          \ans{$\det(B) = 3(2) - 6(1) = 6 - 6 = 0$}

    \item What does the value of $\det(B)$ tell you geometrically about the two column
          vectors of $B$? Explain in a complete sentence.\\[4pt]
          \ansline{Since $\det(B) = 0$, the two column vectors $\begin{bmatrix}3\\1\end{bmatrix}$
          and $\begin{bmatrix}6\\2\end{bmatrix}$ are parallel (one is a scalar multiple of the other),}
          \ansline{so they span a line rather than a parallelogram; the area of the ``parallelogram'' is 0.}

    \item Can $B$ be inverted? Justify your answer.\\[4pt]
          \ansline{No. Since $\det(B) = 0$, the matrix $B$ is singular and has no inverse.}
\end{enumerate}
\end{tcolorbox}

\end{document}
